% Week 5 Session A Lab 1 — Rainfall forecasting & alert system
% Build: pdflatex Week5_SessionA_Lab1_Report.tex from Lab-reports/ (twice)
% Screenshots:
%   week5_session_a_lab1_terminal.png
%   week5_session_a_lab1_streamlit.png
% Code: week5_session_a_lab1_files/
%
\documentclass[11pt,a4paper]{article}

\usepackage[margin=2.5cm]{geometry}
\usepackage{graphicx}
\newcommand{\termfig}[1]{%
  \includegraphics[width=\linewidth,height=0.72\textheight,keepaspectratio]{#1}%
}
\newcommand{\streamlitfig}[1]{%
  \includegraphics[width=\linewidth,height=0.58\textheight,keepaspectratio]{#1}%
}
\usepackage{booktabs}
\usepackage{xurl}
\usepackage{hyperref}
\usepackage{parskip}
\usepackage{enumitem}
\usepackage{xcolor}
\usepackage{fvextra}
\usepackage[most]{tcolorbox}

\newcommand{\studentname}{Mahmudul Hasan}
\newcommand{\studentid}{4125999049}
\newcommand{\reportdate}{May 21, 2026}

\makeatletter
\g@addto@macro\UrlBreaks{\do\ }
\makeatother

\newcommand{\LabCodeRoot}{week5_session_a_lab1_files}

\newcommand{\filebox}[2]{%
\begin{tcolorbox}[
  enhanced,
  colback=gray!6,
  colframe=black!55,
  title={#1},
  fonttitle=\bfseries\ttfamily\footnotesize,
  arc=1.5mm,
  boxrule=0.7pt,
  breakable,
  width=\linewidth,
  before skip=0.8em,
  after skip=0.8em,
]
\VerbatimInput[fontsize=\footnotesize,breaklines=true,breakanywhere=true]{\LabCodeRoot/#2}
\end{tcolorbox}%
}

\title{Lab Report: Week 5 Session A Lab 1\\Rainfall Forecasting \& Alert System}
\author{\studentname \quad (Student ID: \studentid)}
\date{\reportdate}

\begin{document}
\maketitle

\begin{abstract}
This report documents Week~5 Session~A Lab~1: a \textbf{rainfall forecasting and alert system} using the OpenWeather \textbf{Current Weather API} (\texttt{2.5/weather}), threshold-based alerts (GREEN / YELLOW / RED in mm/h), and a Streamlit monitoring dashboard. Work was completed in \texttt{week5\_session\_a\_lab1/} with API keys stored only in the environment variable \texttt{OPENWEATHER\_API\_KEY}. After initial \texttt{401 Unauthorized} (key activation delay), five cities were fetched successfully (0.00~mm/h, GREEN). Invalid-key handling was verified with \texttt{OPENWEATHER\_API\_KEY=bad}. Evidence: terminal screenshot (401, multi-city runs, alert history CSV) and Streamlit screenshot; full source in Appendix~\ref{app:weather}--\ref{app:promptlog}.
\end{abstract}

\section{Introduction}
Hydrological monitoring often combines \textbf{external weather APIs}, \textbf{domain thresholds}, and a \textbf{user-facing dashboard}. This lab practices integrating those pieces with AI-assisted development while validating rainfall semantics: OpenWeather may expose \texttt{rain.1h}, \texttt{rain.3h}, or no rain field when conditions are dry---missing rain must be treated as 0~mm/h for consistent alerting. The implementation builds on prior course work (Week~4 integration patterns) but uses city-based OpenWeather queries as specified in the lab guide.

\section{Objectives}
\begin{itemize}[noitemsep]
  \item Integrate OpenWeather API for current rainfall by city.
  \item Implement GREEN / YELLOW / RED threshold alerting with history.
  \item Build a Streamlit dashboard with refresh and alert display.
  \item Validate behavior including API failure and dry conditions.
\end{itemize}

\section{Methods}

\subsection{API and rainfall extraction (Exercise 1)}
\textbf{Endpoint:} \texttt{https://api.openweathermap.org/data/2.5/weather?q=\{city\}\&units=metric\&appid=KEY}.

\textbf{Rain field:} Prefer \texttt{rain.1h} (mm in last hour); fallback \texttt{rain.3h}; if absent, rainfall = 0~mm/h (dry).

\textbf{Client features:} \texttt{openweather\_client.py} --- \texttt{fetch\_current\_weather()}, in-memory cache (TTL from \texttt{CACHE\_TTL\_SEC}), logging, errors for 401 (invalid key), 404 (city not found), network/JSON failures.

\textbf{Configuration:} \texttt{config.py} --- \texttt{load\_settings()} reads \texttt{OPENWEATHER\_API\_KEY} (required), default city, thresholds, cache TTL. No secrets in source files.

\subsection{Alert logic (Exercise 2)}
\texttt{alerts.py} --- \texttt{check\_alert(rainfall\_mm\_h, settings)}:
\begin{itemize}[noitemsep]
  \item GREEN: rainfall $< 10$~mm/h
  \item YELLOW: $10 \le$ rainfall $\le 20$~mm/h
  \item RED: rainfall $> 20$~mm/h
\end{itemize}
\texttt{append\_alert\_history()} writes to \texttt{data/alert\_history.csv}. \texttt{main.py} orchestrates fetch $\rightarrow$ alert $\rightarrow$ history.

\subsection{Dashboard (Exercise 3)}
\texttt{dashboard.py} --- Streamlit UI: city input, Refresh (bypasses cache on refresh), large rainfall metric, color-coded alert, last update, history table from CSV. Launched with \texttt{python3 -m streamlit run dashboard.py}.

\subsection{Validation (Exercise 4)}
Dry multi-city CLI runs; invalid API key (401); CSV history inspection; dashboard at \texttt{http://localhost:8501}.

\section{Results}

\subsection{File inventory}
\begin{table}[htbp]
  \centering
  \caption{Submitted files in \texttt{week5\_session\_a\_lab1/}.}
  \label{tab:files}
  \small
  \begin{tabular}{@{}ll@{}}
    \toprule
    File & Role \\
    \midrule
    \texttt{openweather\_client.py} & API + cache + rain extraction \\
    \texttt{alerts.py} & Threshold alerts + CSV history \\
    \texttt{config.py} & Environment settings \\
    \texttt{main.py} & CLI pipeline \\
    \texttt{dashboard.py} & Streamlit monitor \\
    \texttt{data/alert\_history.csv} & Alert log \\
    \texttt{prompt\_*.txt}, \texttt{prompt\_log.md} & Documentation \\
    \bottomrule
  \end{tabular}
\end{table}

\subsection{Terminal verification}
Figure~\ref{fig:terminal} shows the full CLI session: initial \texttt{401 Invalid API key} (key not yet active), successful fetches after activation, five cities, and \texttt{alert\_history.csv}.

\begin{figure}[htbp]
  \centering
  \termfig{week5_session_a_lab1_terminal.png}
  \caption{Terminal: 401 before key activation; successful runs (Dhaka, Beijing, Delhi, Moscow, Islamabad); alert history CSV.}
  \label{fig:terminal}
\end{figure}

Representative successful output:
\begin{verbatim}
Dhaka,BD: rain=0.00 mm/h (few clouds) -> GREEN
\end{verbatim}

Additional test:
\begin{verbatim}
OPENWEATHER_API_KEY=bad python3 main.py Dhaka,BD
API ERROR: Invalid API key (401 Unauthorized)
\end{verbatim}

\subsection{Alert history}
After five CLI runs, \texttt{alert\_history.csv} contained five rows (all GREEN, 0.00~mm/h) with UTC timestamps and city names---confirming Exercise~2 persistence.

\subsection{Streamlit dashboard}
Figure~\ref{fig:streamlit} shows the monitoring UI at \texttt{localhost:8501}: current rainfall, alert status, and history aligned with CSV data.

\begin{figure}[htbp]
  \centering
  \streamlitfig{week5_session_a_lab1_streamlit.png}
  \caption{Streamlit dashboard: city selector, rainfall metric, color-coded alert, history table.}
  \label{fig:streamlit}
\end{figure}

Table~\ref{tab:validation} maps lab checklist items to outcomes.

\begin{table}[htbp]
  \centering
  \caption{Exercise 4 validation checklist.}
  \label{tab:validation}
  \begin{tabular}{@{}ll@{}}
    \toprule
    Check & Outcome \\
    \midrule
    Dry city $\rightarrow$ GREEN & Five cities at 0.00~mm/h, GREEN \\
    Invalid API key & 401 user-visible (before activation and with \texttt{bad} key) \\
    CSV append & Five history rows \\
    Dashboard & Runs; reflects stored history \\
    \bottomrule
  \end{tabular}
\end{table}

\section{Analysis and validation}

\textbf{Domain note:} All test cities returned dry conditions (no \texttt{rain} object), correctly mapped to 0~mm/h and GREEN. YELLOW/RED would appear when \texttt{rain.1h} exceeds 10 and 20~mm/h---logic is in \texttt{check\_alert()} and can be tested with instructor fixtures or rainy-period cities.

\textbf{API key lifecycle:} OpenWeather free keys often return 401 until activated; documenting 401 $\rightarrow$ success in one terminal screenshot satisfies both error-handling and operational verification.

\textbf{Security:} Key supplied only via \texttt{export OPENWEATHER\_API\_KEY}; not committed to repository or report.

\textbf{Comparison to Week~4B:} Week~4 used Open-Meteo without a key; Week~5 Lab~1 requires OpenWeather city queries and explicit mm/h thresholds for operational alerting.

\section{Discussion}
\begin{enumerate}[noitemsep]
  \item Treating missing \texttt{rain} as zero is essential---otherwise dry weather would be undefined.
  \item A single terminal screenshot spanning 401 and success is efficient evidence for graders.
  \item Streamlit reads CSV only---appropriate separation from live API except on Refresh.
  \item Optional OpenCode prompts (\texttt{prompt\_weather\_api.txt}, etc.) are available; baseline code ran without further AI edits.
\end{enumerate}

\section{Problems encountered}
\texttt{streamlit: command not found} on first launch---resolved with \texttt{python3 -m streamlit run}. New API key returned 401 until activation (~15+ minutes). No YELLOW/RED cases during testing because all sampled cities were dry.

\section{Lessons learned}
Document the exact JSON rain field and mm/h units in \texttt{prompt\_log.md}. Test invalid keys explicitly. Use environment variables for secrets. Plan time for OpenWeather key activation before the lab deadline.

\section{Conclusion}
Week~5 Session~A Lab~1 objectives were met: OpenWeather integration, threshold alerts with CSV history, Streamlit dashboard, and validation including 401 handling and multi-city dry runs. Submitted artifacts: \texttt{week5\_session\_a\_lab1/}, \texttt{prompt\_log.md}, terminal and Streamlit screenshots.

\appendix

\section{Weather client: \texttt{openweather\_client.py}}
\label{app:weather}
\filebox{\texttt{openweather\_client.py} (full file)}{openweather_client.py}

\section{Alerts: \texttt{alerts.py}}
\label{app:alerts}
\filebox{\texttt{alerts.py} (full file)}{alerts.py}

\section{CLI: \texttt{main.py}}
\label{app:main}
\filebox{\texttt{main.py} (full file)}{main.py}

\section{Dashboard: \texttt{dashboard.py}}
\label{app:dashboard}
\filebox{\texttt{dashboard.py} (full file)}{dashboard.py}

\section{Configuration: \texttt{config.py}}
\label{app:config}
\filebox{\texttt{config.py} (full file)}{config.py}

\section{Submitted prompt log: \texttt{prompt\_log.md}}
\label{app:promptlog}
\filebox{\texttt{prompt\_log.md} (full file)}{prompt_log.md}

\end{document}
